\chapter{Viraje coordinado y factores de carga}\label{ch:turn}
\objetivos{Derivar radio, razón de giro y carga; incluir pérdida y resistencia; separar viraje sostenido e instantáneo y trayectoria aérea y terrestre.}

\section{Balance geométrico del viraje nivelado}
Se considera vuelo coordinado, sin resbalamiento significativo, altura y rapidez constantes, viento nulo y empuje paralelo a la velocidad. Si $\phi$ es alabeo,
\begin{equation}
L\cos\phi=W,\qquad L\sin\phi=mV^2/R.
\end{equation}
La primera ecuación sostiene el peso y la segunda proporciona aceleración centrípeta. Al dividir,
\begin{equation}
n=\frac LW=\frac1{\cos\phi},\qquad
R=\frac{V^2}{g_0\tan|\phi|},\qquad
\dot\chi=\frac{g_0\tan\phi}{V}.
\label{eq:turn}
\end{equation}
El signo de la razón de giro depende del signo de alabeo; el radio se expresa positivo. En $\phi=0$ el radio tiende a infinito. En el límite $|\phi|\to90^\circ$, la carga requerida diverge: la fórmula no permite mantener un viraje nivelado con sustentación finita.

\ilustra{fuerzas-viraje}{Vista frontal esquemática de un viraje: la componente vertical de la sustentación sostiene el peso y la horizontal produce la aceleración centrípeta.}

\section{Carga y pérdida en maniobra}
Con $L=nW$,
\begin{equation}
C_L=\frac{2nW}{\rho V^2S},\qquad
V_{\rm stall}(n)=\sqrt n\,V_{\rm stall}(1).
\end{equation}
Un viraje de $60^\circ$ exige $n=2$ y eleva la velocidad de pérdida por $\sqrt2$, si $C_{L\max}$ permanece igual. La masa y densidad modifican TAS de pérdida; su EAS depende de masa y factor de carga. Una maniobra transitoria puede introducir otros efectos de pérdida y cargas que no recoge el equilibrio estacionario.

\ilustra{carga-viraje}{Factor de carga y multiplicador de velocidad de pérdida en función del alabeo. La figura no incorpora un límite estructural específico de una aeronave real.}

\section{Resistencia y empuje en viraje}
La polar da
\begin{equation}
D(V,n)=AV^2+\frac{Bn^2}{V^2},
\label{eq:dragturn}
\end{equation}
con los $A,B$ de vuelo nivelado para $n=1$. La parte cuadrática asociada a sustentación se multiplica por $n^2$. A igual velocidad, doblar carga cuadruplica esa parte de la resistencia, aunque la parte parásita no cambie. Por eso mantener velocidad y altura requiere más empuje o potencia.

Con empuje disponible $T_d$ independiente de $n$, el límite de carga sostenida por propulsión es
\begin{equation}
n_T^2(V)=\frac{[T_d(V)-AV^2]V^2}{B},
\end{equation}
si el miembro derecho es positivo. Se combina con
\begin{equation}
n_L(V)=\frac{\rho V^2SC_{L\max}}{2W},\qquad n\leq n_{\rm estructural}.
\end{equation}
La carga admisible es la menor de las fronteras aplicables. Un equilibrio de fuerza normal que necesita más empuje del disponible puede alcanzarse transitoriamente perdiendo velocidad o altura, pero no sostenerse bajo las mismas condiciones.

\section{Radio mínimo y razón máxima}
En función de carga,
\begin{equation}
R=\frac{V^2}{g_0\sqrt{n^2-1}},\qquad
|\dot\chi|=\frac{g_0\sqrt{n^2-1}}{V},\quad n>1.
\end{equation}
Reducir radio favorece baja velocidad y alta carga; aumentar razón de giro favorece alta carga respecto a velocidad. No se pueden variar independientemente porque sustentación, estructura y motor acoplan $n$ y $V$.

Para encontrar un óptimo sostenido se recorre la velocidad admisible, se calcula $n_{\max}(V)$ con todas las restricciones y se evalúan radio y razón. Se incluyen los puntos de intersección de fronteras: un óptimo suele aparecer en una esquina de la envolvente y no donde una función suave tiene derivada cero.

\ilustra{radio-viraje}{Radio de giro en calma para dos alabeos fijos. Cada punto sigue necesitando sustentación y empuje suficientes; la figura geométrica no constituye por sí sola una envolvente.}

\section{Viraje instantáneo y sostenido}
El viraje instantáneo usa la capacidad aerodinámica y estructural durante un intervalo corto, admitiendo pérdida de energía. El sostenido exige que el motor compense la resistencia y que altura y velocidad permanezcan constantes. No es correcto presentar un radio calculado con $C_{L\max}$ como radio sostenido sin comprobar potencia.

La altura de energía vuelve a dar una prueba útil: si $TV<DV$ y la altura no cambia, la velocidad debe caer. Un cálculo que mantiene simultáneamente velocidad, altura y empuje insuficiente viola el balance energético aunque el radio geométrico esté bien calculado.

\caso{Ejemplo original: $35^\circ$ a $90\SIu{m\,s^{-1}}$}{Para $\phi=35^\circ$, $n=1.2208$. A $V=90\SIu{m\,s^{-1}}$, $R\simeq1.180\SIu{km}$ y $\dot\chi\simeq0.07630\SIu{rad\,s^{-1}}=4.37^\circ/\mathrm s$. La fuerza horizontal por unidad de masa es $ng_0\sin\phi$, que coincide con $V^2/R$. Para un avión concreto se calcula después $C_L$ y $D(V,n)$ y se comparan con sus límites; el dato geométrico no sustituye esas comprobaciones.}

\section{Efecto del viento}
En viento uniforme, el círculo respecto al aire se traslada con la masa de aire. La trayectoria respecto al terreno deja de ser, en general, un círculo. Con velocidad angular aérea $\omega$ constante y viento horizontal $(U_N,U_E)$,
\begin{align}
x_N(t)&=x_{N0}+\frac V\omega[\sin(\chi_0+\omega t)-\sin\chi_0]+U_Nt,\\
x_E(t)&=x_{E0}-\frac V\omega[\cos(\chi_0+\omega t)-\cos\chi_0]+U_Et.
\end{align}
Tras una vuelta aérea completa, el desplazamiento terrestre es el viento multiplicado por el periodo. Mantener un círculo sobre el terreno requiere variar rumbo, alabeo o velocidad y resolver otro problema.

\section{Conexión con seis grados de libertad}
Un viraje coordinado ideal de punto material no asegura equilibrio de todos los momentos. Las tasas cuerpo se obtienen de la cinemática y pueden generar términos giroscópicos. El timón, el diedro y los mandos deben mantener el resbalamiento requerido. En un simulador se comprobarán $\beta$, tasas, altura y velocidad, además de la forma de la trayectoria.

Para comparar actuaciones entre alturas se debe precisar si se conserva TAS o EAS. A igual EAS, la presión dinámica es igual, pero TAS y radio para un mismo alabeo aumentan cuando disminuye densidad. A igual TAS, la sustentación adimensional necesaria aumenta con altura y puede aparecer pérdida.

\section{Taller resuelto: Radio mínimo sostenido con tres fronteras}

\noindent\textbf{Nivel:} Intermedio.

\apartado{1. Enunciado}{Busque por barrido el radio mínimo sostenido entre pérdida y TAS 180 m/s para masa 4500 kg, rho=1.0, S=30 m², polar 0.024+0.046CL², CLmax=1.6, carga estructural 3 y empuje 7200 N.}

\apartado{2. Datos}{Barrido de 20001 velocidades; si n<=1 no existe viraje nivelado de radio finito.}

\apartado{3. Incógnitas}{Las magnitudes y comprobaciones solicitadas en el caso.}

\apartado{4. Sistema de referencia}{Viraje coordinado nivelado en calma; TAS.}

\apartado{5. Hipótesis físicas}{Empuje constante, polar fija, sin resistencia adicional de mandos ni límite de Mach.}

\apartado{6. Modelo matemático}{Radio mínimo sostenido con tres fronteras}

\apartado{7. Ecuaciones simbólicas}{Se parte de los balances y definiciones siguientes.}

\begin{equation}n_T=\sqrt{(T-AV^2)V^2/B}\end{equation}

\begin{equation}n_L=\rho V^2SC_{Lmax}/(2W)\end{equation}

\begin{equation}n_{max}(V)=\min(n_T,n_L,3)\end{equation}

\begin{equation}R(V)=V^2/[g_0\sqrt{n_{max}^2-1}]\end{equation}

\apartado{8. Desarrollo algebraico}{Las operaciones se realizan en este orden, conservando las condiciones del modelo.}

\begin{enumerate}

\item Se recorre sólo el intervalo de velocidad declarado.

\item Para cada punto se toma el menor de los tres límites de carga.

\item Se descartan puntos sin carga superior a uno y se minimiza el radio.

\item Se revisa la frontera activa cerca del mínimo y se afina la rejilla si se requiere más precisión.

\end{enumerate}

\apartado{9. Sustitución numérica}{Se usan las magnitudes del apartado de datos en SI; los ángulos se convierten a radianes antes de evaluar funciones trigonométricas. Las densidades indicadas como ISA se obtienen de la capa correspondiente.}

\noindent TAS de rejilla óptima:\quad $58.18696160080154\simeq 58.18696\,\mathrm{m\,s^{-1}}$.\par

\noindent Carga en el punto:\quad $1.841320600658072\simeq 1.841321\,$.\par

\noindent Radio de rejilla mínimo:\quad $223.3005549791824\simeq 223.3006\,\mathrm m$.\par

\apartado{10. Resultados con unidades}{Los resultados pertenecen al modelo y al intervalo de datos declarados.}

\begin{center}\begin{tabular}{lr}\toprule Magnitud&Resultado\\\midrule

TAS de rejilla óptima & $58.18696\,\mathrm{m\,s^{-1}}$\\

Carga en el punto & $1.841321\,$\\

Radio de rejilla mínimo & $223.3006\,\mathrm m$\\

Límite motor en el punto & $1.841413\,$\\

Límite sustentación en el punto & $1.841321\,$\\

\bottomrule\end{tabular}\end{center}

\apartado{11. Interpretación}{El óptimo de rejilla puede estar junto a una intersección de límites; no debe presentarse como raíz suave universal.}

\apartado{12. Comprobación independiente y alcance}{Refinar la rejilla y resolver las intersecciones de las fronteras próximas; comprobar D<=T, CL<=CLmax y n<=3.}

\apartado{13. Limitaciones}{La resolución de velocidad limita la precisión del óptimo y no se oculta con muchas cifras impresas.}

\par\smallskip\noindent{\footnotesize\textit{Datos y solución originales. Las sustituciones registradas se reproducen con SymPy y Studium; eso no establece validación experimental ni revisión experta.}}


\repaso{¿Mantiene el mismo radio un avión que conserva EAS y alabeo al subir?}{No. Si la densidad baja, la TAS aumenta y $R\propto V_{\rm TAS}^2$. El equilibrio aerodinámico puede conservar un $C_L$ parecido a masa fija, pero la geometría de la trayectoria cambia.}
\resumen{El viraje nivelado necesita sustentación superior al peso. Radio y razón de giro son relaciones geométricas sujetas a límites de sustentación, estructura y propulsión. La trayectoria terrestre cambia con viento, y el trim completo requiere también equilibrio de momentos.}
\fuentes{Guía ULE, viraje coordinado y actuaciones de aviones de reacción y hélice. Deducción propia desde balances vertical y centrípeto; ecuaciones verificadas dimensionalmente y en el caso numérico original.}
